\documentclass[12pt,a4paper]{article}

\usepackage[spanish,es-tabla]{babel}
\usepackage[utf8]{inputenc}
\usepackage[T1]{fontenc}
\usepackage{amsmath}
\usepackage{siunitx}
\usepackage{booktabs}
\usepackage{geometry}
\usepackage{tikz}
\usepackage{hyperref}
\usepackage[numbers,sort&compress]{natbib}
\usetikzlibrary{arrows.meta,babel,positioning}

\geometry{margin=2.5cm}
\sisetup{per-mode=symbol,detect-all,separate-uncertainty=true}
\hypersetup{colorlinks=true,linkcolor=blue,citecolor=blue,urlcolor=blue}

\title{Caracterización de una metasuperficie electromagnética reconfigurable en microondas mediante parámetros $S$}
\author{Propuesta experimental conceptual de laboratorio}
\date{\today}

\begin{document}

\maketitle

\begin{abstract}
Se presenta una propuesta experimental conceptual, no destinada a ser ejecutada durante el curso, para caracterizar una metasuperficie electromagnética reconfigurable en el rango de microondas mediante la medición de parámetros $S$. Se propone medir los coeficientes complejos de reflexión $S_{11}$ y transmisión $S_{21}$ con un analizador vectorial de redes, y a partir de ellos derivar reflectancia $R$, transmitancia $T$, absorción $A$, frecuencia de resonancia $f_0$ y factor de calidad $Q$. La propuesta organiza el diseño del montaje, el fundamento físico, los datos esperados, la estimación de incertidumbres y la utilidad tecnológica de este tipo de superficies artificiales.
\end{abstract}

\section{Título específico y pregunta experimental}

\textbf{Caracterización de una metasuperficie electromagnética reconfigurable en microondas mediante parámetros $S$: control de reflexión, transmisión y absorción resonante.}

La pregunta experimental propuesta es: ¿cómo cambiarían la reflexión, la transmisión y la absorción de una onda de microondas al iluminar una metasuperficie reconfigurable y variar su condición de resonancia? En lugar de reportar datos medidos, este documento formula un experimento físicamente consistente para obtener $S_{11}(f)$ y $S_{21}(f)$, transformar esas magnitudes en observables energéticos y discutir qué comportamiento se esperaría observar cerca de una resonancia electromagnética.

\section{Alcance y carácter propositivo de la propuesta}

Esta es una propuesta experimental conceptual, no destinada a ser ejecutada durante el curso. Su alcance no es construir, calibrar ni operar materialmente el sistema, sino formular un experimento posible, coherente con la electrodinámica clásica y evaluable por sus variables, procedimientos y criterios de análisis.

El objetivo académico es diseñar un montaje propuesto, una base teórica, una estructura de datos esperados, un análisis de incertidumbres y una discusión de aplicaciones. La propuesta se organiza en dos niveles de recursos:

\begin{enumerate}
  \item \textbf{Recursos moderados:} analizador vectorial de redes (VNA), dos antenas de microondas, soportes, cables coaxiales, calibración básica y una metasuperficie fabricada como placa de circuito impreso (PCB) con elementos resonantes y, si se desea reconfigurabilidad, diodos varactores o interruptores RF.
  \item \textbf{Recursos avanzados:} cámara anecoica, posicionadores angulares, fabricación de precisión, control de polarización, simulación electromagnética de onda completa y procedimientos de calibración de espacio libre más rigurosos.
\end{enumerate}

El diseño conecta con contenidos centrales del curso: condiciones de frontera, propagación en medios, impedancia, reflexión, transmisión, dispersión, resonancia, pérdidas y scattering. Así, el énfasis está en traducir conceptos de electrodinámica en magnitudes observables y no en afirmar una ejecución que no forma parte del alcance.

\section{Utilidad tecnológica y aplicaciones}

Una metasuperficie es una estructura artificial delgada capaz de modificar la respuesta de una onda electromagnética mediante subunidades resonantes menores que la longitud de onda. Su utilidad tecnológica radica en que permite controlar amplitud, fase, polarización, dirección o absorción de ondas sin recurrir necesariamente a dispositivos volumétricos. En ese sentido, conecta la electrodinámica clásica con tecnologías modernas de control de ondas.

Las aplicaciones relevantes incluyen:

\begin{itemize}
  \item absorbedores electromagnéticos de banda estrecha o ajustable;
  \item cámaras anecoicas y terminaciones de baja reflexión;
  \item blindaje y compatibilidad electromagnética;
  \item telecomunicaciones y superficies inteligentes reconfigurables;
  \item control de haz y redirección de frentes de onda;
  \item sensores basados en desplazamiento de resonancia;
  \item radar y control de sección eficaz;
  \item formación de imágenes en microondas;
  \item dispositivos de terahercios;
  \item control de polarización;
  \item reducción de reflexiones no deseadas en interfaces.
\end{itemize}

Aunque la propuesta no se ejecute durante el curso, su valor está en traducir conceptos de electrodinámica en magnitudes medibles: $S_{11}$, $S_{21}$, $R$, $T$, $A$, $f_0$, $Q$ e impedancia efectiva.

\section{Objetivos generales y específicos}

\subsection{Objetivo general}

Proponer y diseñar un experimento de caracterización electromagnética de una metasuperficie reconfigurable en microondas mediante parámetros $S$, para analizar reflexión, transmisión, absorción resonante, frecuencia de resonancia, factor de calidad e impedancia efectiva.

\subsection{Objetivos específicos}

\begin{itemize}
  \item Formular el diseño conceptual de una metasuperficie resonante reconfigurable en el rango de microondas.
  \item Proponer un montaje de medición con VNA, antenas y control de geometría de incidencia.
  \item Definir los observables complejos $S_{11}$ y $S_{21}$ como base para estimar $R$, $T$ y $A$.
  \item Establecer un procedimiento reproducible aunque el montaje propuesto no se ejecute durante el curso.
  \item Proponer un análisis de incertidumbres asociado a calibración, alineación, distancia, polarización, ruido y repetibilidad.
  \item Discutir aplicaciones tecnológicas vinculadas con absorción, blindaje, telecomunicaciones, sensores, radar y control de polarización.
\end{itemize}

\section{Fundamento teórico}

\subsection{Parámetros $S$ en microondas}

En microondas es más natural caracterizar un dispositivo por ondas incidentes y salientes que por voltajes y corrientes locales. Para una muestra iluminada desde el puerto 1, $S_{11}$ representa el coeficiente complejo de reflexión y $S_{21}$ representa el coeficiente complejo de transmisión hacia el puerto 2. Ambos dependen de la frecuencia y contienen módulo y fase:
\begin{equation}
  S_{ij}(f)=|S_{ij}(f)|e^{i\phi_{ij}(f)}.
\end{equation}

Si las antenas, líneas y puertos están normalizados a la misma impedancia característica, las fracciones de potencia reflejada y transmitida pueden aproximarse como
\begin{equation}
  R(f)=|S_{11}(f)|^2,
  \qquad
  T(f)=|S_{21}(f)|^2.
\end{equation}
La absorción esperada se infiere por balance energético:
\begin{equation}
  A(f)=1-R(f)-T(f).
  \label{eq:absorcion}
\end{equation}
Esta relación es válida como modelo operativo cuando las pérdidas laterales, la difracción fuera del cono de medición y las reflexiones del entorno están controladas o incorporadas como incertidumbres sistemáticas.

\subsection{Impedancia efectiva y condición de absorción}

Una metasuperficie absorbente puede interpretarse como una frontera efectiva con impedancia superficial ajustable. Para incidencia normal, la reflexión se minimiza cuando la impedancia efectiva de la superficie se aproxima a la impedancia del espacio libre $Z_0\approx\SI{377}{\ohm}$. Si además existen pérdidas internas, la energía que no se refleja ni se transmite se disipa en la estructura.

Una forma usual de estimar la impedancia efectiva normalizada $z_\mathrm{eff}=Z_\mathrm{eff}/Z_0$ a partir de parámetros $S$ en una celda equivalente es
\begin{equation}
  z_\mathrm{eff}=\pm\sqrt{\frac{(1+S_{11})^2-S_{21}^2}{(1-S_{11})^2-S_{21}^2}},
  \label{eq:impedancia}
\end{equation}
donde el signo se selecciona imponiendo pasividad. En esta propuesta, la ecuación~\eqref{eq:impedancia} sirve como referencia conceptual: el análisis principal se centraría en $S_{11}$, $S_{21}$, $R$, $T$ y $A$.

\subsection{Resonancia, modelo LC y scattering}

Las unidades de una metasuperficie pueden diseñarse como resonadores sublongitud de onda. La geometría metálica aporta una inductancia efectiva $L$, las brechas o separaciones aportan una capacitancia efectiva $C$ y las pérdidas conductivas o dieléctricas se representan mediante una resistencia efectiva. En primera aproximación, la frecuencia de resonancia esperada cumple
\begin{equation}
  f_0=\frac{1}{2\pi\sqrt{LC}}.
\end{equation}

Si la metasuperficie es reconfigurable, por ejemplo mediante diodos varactores, un cambio de capacitancia desplazaría $f_0$. Cerca de la resonancia se esperaría una variación rápida de fase, un mínimo de $|S_{11}|$ si hay adaptación de impedancia, una caída de $|S_{21}|$ si la superficie bloquea la propagación y un máximo de $A$ si las pérdidas están ajustadas para disipar energía. El factor de calidad se definiría como
\begin{equation}
  Q=\frac{f_0}{\Delta f},
\end{equation}
donde $\Delta f$ es el ancho de banda asociado al criterio elegido, por ejemplo el ancho a mitad de altura del pico de absorción esperado.

\section{Materiales y recursos propuestos}

La propuesta distingue entre un montaje de recursos moderados y una versión avanzada. Ambos niveles persiguen la misma lógica física: iluminar la muestra, registrar $S_{11}$ y $S_{21}$, y transformar esos datos esperados en magnitudes energéticas.

\begin{table}[htbp]
  \centering
  \caption{Recursos propuestos para dos niveles de implementación.}
  \begin{tabular}{p{0.23\linewidth}p{0.34\linewidth}p{0.34\linewidth}}
    \toprule
    Elemento & Recursos moderados & Recursos avanzados \\
    \midrule
    Instrumentación & VNA de dos puertos, cables coaxiales y kit de calibración. & VNA de mayor rango dinámico, calibración de espacio libre y verificación trazable. \\
    Iluminación & Dos antenas de bocina o antenas de banda ancha alineadas. & Antenas calibradas, control de polarización y posicionadores angulares. \\
    Muestra & PCB con patrones metálicos resonantes y carga reconfigurable simple. & Fabricación de precisión, sustratos caracterizados y red de polarización RF optimizada. \\
    Entorno & Banco con absorbentes alrededor de la trayectoria principal. & Cámara anecoica y control de reflexiones múltiples. \\
    Modelado & Modelo LC y simulación simplificada. & Simulación de onda completa con barridos paramétricos. \\
    \bottomrule
  \end{tabular}
\end{table}

\section{Montaje propuesto y procedimiento de medición}

El montaje propuesto ubicaría la metasuperficie entre dos antenas conectadas al VNA. La antena emisora definiría la onda incidente y la antena receptora registraría la onda transmitida. La reflexión se mediría desde el puerto de entrada, idealmente con calibración que desplace el plano de referencia hasta la posición de la muestra o con una medición de referencia que permita sustraer efectos del espacio libre.

\begin{figure}[htbp]
  \centering
  \begin{tikzpicture}[>=Stealth,scale=0.92]
    \node[draw,rounded corners,minimum width=2.4cm,minimum height=1.1cm] (vna) at (0,0) {VNA};
    \node[draw,minimum width=1.3cm,minimum height=1.0cm,right=2.0cm of vna] (tx) {Antena 1};
    \node[draw,minimum width=0.25cm,minimum height=2.8cm,right=2.0cm of tx,fill=blue!12] (ms) {};
    \node[above=0.1cm of ms] {metasuperficie};
    \node[draw,minimum width=1.3cm,minimum height=1.0cm,right=2.0cm of ms] (rx) {Antena 2};
    \draw[->,thick] (vna.east) -- node[above] {puerto 1} (tx.west);
    \draw[->,thick] (rx.west) -- node[above] {$S_{21}$} (ms.east);
    \draw[->,thick] (tx.east) -- node[above] {onda incidente} (ms.west);
    \draw[->,thick] (ms.east) -- node[above] {transmisión} (rx.west);
    \draw[<-,thick,red!70!black] (tx.east) -- ++(0.7,-0.8) node[right] {$S_{11}$};
    \draw[->,thick] (vna.south) .. controls (4.5,-2.0) and (7.2,-1.9) .. node[below] {puerto 2} (rx.south);
  \end{tikzpicture}
  \caption{Esquema del montaje propuesto para caracterización en espacio libre de una metasuperficie con VNA y dos antenas.}
  \label{fig:montaje}
\end{figure}

El procedimiento reproducible propuesto sería:

\begin{enumerate}
  \item Definir el rango de frecuencias de interés, por ejemplo una banda de microondas alrededor de la resonancia diseñada.
  \item Calibrar el VNA y registrar una medición de referencia sin muestra para estimar pérdidas de trayecto y reflexiones residuales.
  \item Alinear antenas y muestra para incidencia normal, con polarización compatible con la geometría resonante.
  \item Ubicar absorbentes alrededor del trayecto directo para reducir reflexiones espurias.
  \item Registrar $S_{11}(f)$ y $S_{21}(f)$ para cada estado de reconfiguración propuesto.
  \item Repetir barridos para estimar repetibilidad y variabilidad por recolocación de la muestra.
  \item Calcular $R(f)$, $T(f)$ y $A(f)$ mediante la ecuación~\eqref{eq:absorcion}.
  \item Identificar $f_0$, $\Delta f$, $Q$ y el desplazamiento de resonancia entre estados.
\end{enumerate}

\section{Datos esperados y visualizaciones previstas}

No se reportan datos medidos. Los datos esperados serían curvas complejas de $S_{11}$ y $S_{21}$ en función de la frecuencia y, si la muestra es reconfigurable, para varios estados de polarización eléctrica, capacitancia o conmutación.

\begin{table}[htbp]
  \centering
  \caption{Estructura esperada del registro de datos propuesto.}
  \begin{tabular}{p{0.18\linewidth}p{0.23\linewidth}p{0.23\linewidth}p{0.25\linewidth}}
    \toprule
    Variable & Magnitud registrada & Magnitud derivada & Comportamiento esperado \\
    \midrule
    Frecuencia $f$ & Barrido del VNA & Eje común de análisis & Resolución suficiente alrededor de $f_0$. \\
    $S_{11}$ & Módulo y fase & $R=|S_{11}|^2$ & Mínimo cerca de adaptación de impedancia. \\
    $S_{21}$ & Módulo y fase & $T=|S_{21}|^2$ & Reducción cerca de una banda resonante. \\
    Estado de control & Sesgo, conmutación o configuración & Desplazamiento de $f_0$ & Sintonización de la respuesta espectral. \\
    $A$ & Derivada de balance energético & $A=1-R-T$ & Pico de absorción si las pérdidas están ajustadas. \\
    \bottomrule
  \end{tabular}
\end{table}

Las gráficas previstas incluirían $|S_{11}|$ y $|S_{21}|$ en escala lineal o dB, fase de $S_{11}$ y $S_{21}$, curvas $R(f)$, $T(f)$ y $A(f)$, comparación de estados reconfigurables, extracción de $f_0$ y estimación de $Q$. Se esperaría observar que la absorción máxima ocurre cuando se combinan baja reflexión, baja transmisión y pérdidas internas no despreciables.

\section{Análisis de incertidumbres y errores}

La incertidumbre dominante esperada provendría de la calibración del VNA, las pérdidas de cables, la alineación angular, la distancia entre antenas y muestra, la polarización incidente, reflexiones múltiples, difracción en los bordes de la muestra y repetibilidad al reposicionar el arreglo. En un entorno moderado, estas fuentes pueden ser comparables al cambio de respuesta que se intenta observar, por lo que el diseño debe incluir mediciones de referencia y repeticiones.

Para las magnitudes derivadas se propagarían incertidumbres desde los módulos de $S_{11}$ y $S_{21}$. En forma diferencial,
\begin{equation}
  \sigma_R\approx 2|S_{11}|\sigma_{|S_{11}|},
  \qquad
  \sigma_T\approx 2|S_{21}|\sigma_{|S_{21}|},
\end{equation}
y, si $R$ y $T$ se tratan como independientes en primera aproximación,
\begin{equation}
  \sigma_A\approx\sqrt{\sigma_R^2+\sigma_T^2}.
\end{equation}
También se estimaría una incertidumbre en $f_0$ asociada a la resolución de frecuencia, al ruido en la curva y al criterio usado para localizar el máximo o mínimo resonante. La incertidumbre de $Q$ dependería de la extracción de $\Delta f$ y del suavizado aplicado a los datos esperados.

\section{Seguridad y buenas prácticas}

El montaje propuesto evita altas potencias; aun así, un sistema de microondas requiere buenas prácticas. La potencia de salida del VNA debe mantenerse en niveles seguros para los equipos y para las personas, las antenas no deben apuntarse innecesariamente hacia operadores, los conectores RF deben manipularse sin tensión aplicada y se deben respetar los límites de polarización de cualquier elemento reconfigurable.

Si se emplean diodos, fuentes de sesgo o conmutadores, el circuito de polarización debe aislar la señal RF y limitar corrientes. En una versión avanzada, la cámara anecoica y los posicionadores deben operarse siguiendo protocolos institucionales. Toda documentación de montaje debería registrar configuración, geometría, potencia, rango de frecuencias y estado de control de la muestra.

\section{Conclusiones esperadas}

La propuesta permitiría mostrar cómo una metasuperficie resonante puede analizarse mediante parámetros $S$ y cómo esos parámetros se traducen en magnitudes energéticas: reflexión, transmisión y absorción. Conceptualmente, el máximo de absorción se esperaría cerca de una condición de adaptación de impedancia combinada con pérdidas internas, mientras que la reconfigurabilidad desplazaría la frecuencia de resonancia o modificaría la profundidad de la respuesta.

El valor principal del trabajo no está en afirmar una implementación, sino en construir un puente entre condiciones de frontera, propagación, impedancia, dispersión y scattering. Como propuesta experimental conceptual, el resultado esperado es un diseño claro, reproducible en principio y conectado con aplicaciones actuales de control electromagnético mediante superficies artificiales.

\section{Bibliografía}

\nocite{*}
\bibliographystyle{plainnat}
\bibliography{references}

\end{document}
